\documentclass[conference]{IEEEtran}
\IEEEoverridecommandlockouts
\usepackage{cite}
\usepackage{amsmath,amssymb,amsfonts}
\usepackage{graphicx}
\usepackage{textcomp}
\usepackage{xcolor}
\usepackage{hyperref}
\usepackage{booktabs}

\begin{document}

\title{Memory Game with a Franka Emika Panda:\\
Vision-Based Block Detection and Sequence Replay}

\author{
\IEEEauthorblockN{Ali, Sinan, Izat, Boburjon}
\IEEEauthorblockA{\textit{Robotics and Intelligent Systems Lab 2} \\
\textit{Constructor University Bremen}\\
Bremen, Germany}
}

\maketitle

\begin{abstract}
This paper presents the design and implementation of an interactive
memory game using a Franka Emika Panda robotic arm and an Intel RealSense
RGB-D camera. The robot demonstrates a randomly generated sequence by
pointing at colored blocks on a table, and the player reproduces the
sequence by physically removing blocks from their initial positions. The
system is implemented as a ROS1 package composed of independent nodes for
vision, player selection, game logic, and motion planning, with all
components communicating through standard ROS topics. We describe the
classical computer vision pipeline (HSV thresholding with morphological
operations, depth-based 3D projection, and temporal smoothing), the
slot-based player selection mechanism, and the MoveIt-based motion
planning approach. We discuss three major design decisions made during
development---choosing HSV thresholding over deep learning, pre-detection
image cropping over 3D workspace filtering, and sequence-based motion
dispatch over single-target dispatch---along with the alternatives
considered and the rationale for each choice. The implementation
demonstrates that, with careful tuning and robust state management,
classical computer vision and standard motion planning can produce a
reliable interactive robot system within a semester-length project.
\end{abstract}

\begin{IEEEkeywords}
human-robot interaction, computer vision, motion planning, ROS,
manipulator robotics
\end{IEEEkeywords}

\section{Introduction}

Human-robot interaction (HRI) demonstrations serve a dual purpose in
robotics education: they expose students to the breadth of subsystems
required to build a working robot application, and they produce engaging
artifacts that communicate robotics capabilities to a non-technical
audience. A memory game between a robot and a human is well-suited to
this purpose because it exercises real-time perception, state
management, motion planning, and feedback to the user, while remaining
intuitive enough that an observer immediately understands what the
system is doing.

This paper describes the design and implementation of such a game
system, developed over the course of a semester within the Robotics
and Intelligent Systems Lab 2 course. The hardware platform consists
of a Franka Emika Panda 7-DOF arm and an Intel RealSense RGB-D camera
mounted to observe a tabletop workspace (Fig.~\ref{fig:setup}). The
software is implemented as a ROS1 package using standard middleware
components: OpenCV for image processing, MoveIt for motion planning,
and \texttt{tf2} for coordinate frame management.

\begin{figure}[t]
\centering
\fbox{\rule{0pt}{3.5cm}\rule{0.85\columnwidth}{0pt}}
% \includegraphics[width=\columnwidth]{figures/lab_setup.jpg}
\caption{Lab setup: Franka Emika Panda arm, Intel RealSense camera
on a mounted pole, and three colored blocks on the tabletop workspace.}
\label{fig:setup}
\end{figure}

The remainder of the paper is organized as follows. Section~\ref{sec:game}
describes the game as originally envisioned and as ultimately implemented.
Section~\ref{sec:arch} presents the software architecture (Fig.~\ref{fig:arch}).
Section~\ref{sec:methods} discusses the methods used for the major
subsystems (Figs.~\ref{fig:vision_pipeline} and~\ref{fig:state_machine}).
Section~\ref{sec:decisions} analyzes three significant design decisions
in detail. Sections~\ref{sec:future} and~\ref{sec:lessons} discuss future
work and lessons learned.

\section{Game Description}
\label{sec:game}

This section describes the game from the player's perspective, separated
into the originally envisioned version and the final implementation. The
distinction is important because the differences reflect adaptations
made in response to real-world hardware and environmental constraints
that were not anticipated at the planning stage.

\subsection{Envisioned Game (User Story)}

The originally envisioned interaction is as follows. A player approaches
the table on which four colored blocks (red, green, blue, yellow) are
arranged. The system enters a ready state and indicates that the game
is about to begin. The robot arm autonomously generates a random sequence
of blocks (initially three) and demonstrates it by smoothly moving the
arm above each block in turn, pausing briefly to ``point'' at each one.
The player watches and memorizes the sequence.

After the demonstration, the player reproduces the sequence by interacting
with the blocks in the correct order. The system recognizes each selection
and provides immediate audio-visual feedback: a positive tone for a correct
selection, a negative tone for a wrong selection, and a satisfying
celebration when an entire sequence is completed correctly. The difficulty
escalates: each successful round adds one block to the next sequence, and
the time available for the player to respond decreases over time. The
system maintains a score and persists it across rounds.

The system is robust to environmental factors: it handles changes in
ambient lighting, players wearing colored clothes, and occasional
disturbances such as someone walking past the camera. The interaction
is single-player but designed in a way that would extend naturally to
multi-player tournaments.

\subsection{Final Implementation (User Story)}

The implemented system retains the core game loop of the envisioned
design but operates with several pragmatic constraints. Only three
colors are used in the active game: green, blue, and yellow. Red is
configured but disabled by default due to detection unreliability in
the lab environment, where skin tones, brown wooden surfaces, and the
robot's status indicator overlap with the red HSV range.

The game requires explicit setup: the blocks must be placed on the
table in fixed positions before the game starts, and the player must
remain mostly outside the camera's view during demonstration. Once
the blocks are stable and at least three are detected, the system
records each block's initial position as a ``slot''. The robot then
generates a sequence and performs a hover-point-hover motion above
each block in turn.

The player reproduces the sequence by \emph{physically removing} each
block from its slot in the correct order. The slot-based player
selection mechanism detects a block's absence from its initial
position for longer than a configurable threshold and publishes the
corresponding selection event. After the entire sequence is reproduced,
the game node verifies correctness and updates the score. The current
implementation supports a single round at fixed difficulty; difficulty
progression and multi-player support remain as future work.

A browser-based UI displays the live game state, score, and detected
block list. The UI is primarily for the team and the demonstrator
rather than the player, who interacts directly with the physical
blocks.

\subsection{Differences and Rationale}

The differences between the envisioned and the implemented game fall
into three categories.

\textit{Reduced color palette.} Red was disabled because reliable
discrimination of red against skin tones and warm-colored background
objects requires either narrower HSV ranges (which tend to also reject
real red blocks under varying illumination) or more sophisticated
detection methods than were practical within the project's scope.

\textit{Slot-based instead of gesture-based selection.} The envisioned
design implicitly assumed some form of gesture or hand pose detection.
Slot-based selection (detecting a block's disappearance from its initial
position) was chosen as the simpler and more reliable mechanism. It
relies on the same vision pipeline already used for block detection,
avoids the need for a separate hand-tracking model, and is robust to
the player's skin tone, clothing, and approach direction.

\textit{Fixed difficulty.} The progression mechanic (increasing sequence
length over rounds, decreasing player response time) was deferred to
focus engineering effort on reliable single-round gameplay. Adding
difficulty progression on top of the working baseline is straightforward
in principle but was not prioritized.

\section{System Architecture}
\label{sec:arch}

The software is structured as a ROS1 (Noetic) package named
\texttt{memory\_game}. Functionality is divided across four primary
nodes plus supporting infrastructure (Fig.~\ref{fig:arch}). All
inter-node communication uses standard ROS publish/subscribe topics;
no service or action interfaces are required for the main game loop.

\begin{figure}[t]
\centering
\fbox{\rule{0pt}{5.5cm}\rule{0.85\columnwidth}{0pt}}
% \includegraphics[width=\columnwidth]{figures/architecture.png}
\caption{System architecture: ROS nodes and the topic communication
between them. The vision pipeline produces \texttt{/detected\_blocks},
which the player-selection node and the game node consume. The game
node dispatches sequences to the motion node via
\texttt{/target\_sequence} and receives motion progress via
\texttt{/motion\_status}. The UI server (not shown) subscribes to all
game-relevant topics.}
\label{fig:arch}
\end{figure}

\subsection{ROS Node Organization}

The four primary nodes are:

\begin{itemize}
\item \texttt{vision\_node} (C++): Subscribes to the RGB camera image,
depth image, and camera info topics from the RealSense driver.
Detects colored blocks, projects them into the robot base frame, and
publishes the detection list.
\item \texttt{player\_selection} (Python): Subscribes to the detected
block list. Maintains a record of each block's initial position
(``slot'') and publishes a selection event when a block disappears
from its slot for longer than a threshold time.
\item \texttt{game\_node} (C++): Owns the game state machine. Generates
random sequences, dispatches them to the motion node, receives
selection events from the player\_selection node, and verifies the
player's response.
\item \texttt{motion\_moveit\_node} (C++): Subscribes to target sequence
messages. Executes a hover-point-hover trajectory for each block using
MoveIt-based joint-space planning. Publishes status updates after
each motion stage.
\end{itemize}

A fifth node, the UI server (\texttt{ui/server.py}, Python), provides
a browser-based dashboard. It subscribes to all the game-related
topics, serves a static HTML/JS frontend, and exposes a simple REST
API to start and stop the underlying ROS nodes.

\subsection{Topic Communication}

Table~\ref{tab:topics} summarizes the topics that carry the
game-relevant traffic.

\begin{table}[t]
\centering
\caption{Primary ROS Topics}
\label{tab:topics}
\begin{tabular}{lll}
\toprule
\textbf{Topic} & \textbf{Publisher} & \textbf{Subscribers} \\
\midrule
\texttt{/detected\_blocks} & \texttt{vision\_node} & game, selection, UI \\
\texttt{/player\_selection} & \texttt{player\_selection} & game, UI \\
\texttt{/target\_sequence} & \texttt{game\_node} & motion \\
\texttt{/motion\_status} & \texttt{motion\_moveit\_node} & game, UI \\
\texttt{/game\_state} & \texttt{game\_node} & UI \\
\texttt{/score} & \texttt{game\_node} & UI \\
\bottomrule
\end{tabular}
\end{table}

\subsection{Custom Message Types}

Four custom messages are defined in the \texttt{msg/} directory:

\begin{itemize}
\item \texttt{Block}: a single colored block with integer id, color
name, 3D position, orientation, and a confidence value.
\item \texttt{BlockArray}: a timestamped header and a list of
\texttt{Block} messages, used for the live detection topic.
\item \texttt{BlockSequence}: a timestamped header, a sequence
identifier, and an ordered list of \texttt{Block} messages.
\item \texttt{PlayerSelection}: the id and color of the selected block,
a selection type tag, and a confidence value.
\end{itemize}

The sequence identifier in \texttt{BlockSequence} is also embedded in
all subsequent \texttt{/motion\_status} updates, allowing the game node
to distinguish motion events from the current round from any stragglers
from a previous round.

\subsection{Coordinate Frames}

The camera publishes its outputs in the \texttt{camera\_color\_optical\_frame}.
All block positions are transformed to \texttt{panda\_link0} (the robot
base frame) inside \texttt{vision\_node} using \texttt{tf2}, so that
all downstream consumers operate in a single, robot-relative coordinate
system. The transformation is performed per-detection at the time the
3D position is computed; this requires the TF tree from the camera
frame to the robot base to be valid at the corresponding timestamp.

\section{Methods}
\label{sec:methods}

This section describes the methods used by each of the major subsystems.

\subsection{Vision Pipeline}
\label{sec:vision}

The vision pipeline implemented in \texttt{vision\_node} is a classical
image processing chain (Fig.~\ref{fig:vision_pipeline}). Each color
camera frame triggers the following sequence of operations:

\begin{figure}[t]
\centering
\fbox{\rule{0pt}{6cm}\rule{0.85\columnwidth}{0pt}}
% \includegraphics[width=\columnwidth]{figures/vision_pipeline.png}
\caption{Vision pipeline. A color frame arriving from the RealSense
driver is decoded, synchronized with the closest depth frame, masked
to the configured crop rectangle, converted to HSV, and processed
per-color through thresholding, morphological cleanup, contour
extraction, centroid computation, depth lookup, 3D projection, TF
transformation, and temporal smoothing. Each stage can reject a
candidate, in which case the orchestrator tries the next-best one.}
\label{fig:vision_pipeline}
\end{figure}

\textit{1) Image acquisition and synchronization.} The RGB image is
received via the ROS image topic. Because depth and color are published
independently with their own timestamps, the depth frame is buffered
and the closest-timestamp match (within a configurable tolerance) is
selected for each color frame.

\textit{2) Image-space crop.} Before any color detection, the image
is masked: pixels outside a configurable rectangle (defined by four
pixel-coordinate parameters in YAML) are set to pure black. This crop
ensures that color detection only ever sees the active play area;
objects outside the rectangle---background tables, the robot's own
status LED, players wearing colored clothes---are invisible to all
downstream stages.

\textit{3) Color space conversion.} The cropped BGR image is converted
to HSV using \texttt{cv::cvtColor}. HSV is preferred over BGR because
it separates color identity (hue) from intensity (value) and
saturation, making color matching more robust to lighting variation.

\textit{4) Per-color thresholding.} For each configured color, the
HSV image is thresholded against a configured range using
\texttt{cv::inRange}. Red is configured with two ranges to handle
the wrap-around at hue 0/180; other colors use a single range. The
results are combined via \texttt{cv::bitwise\_or} into a binary mask
per color.

\textit{5) Morphological cleanup.} Each color's mask is processed
with a morphological CLOSE operation (to fill small holes caused by
glare or occlusion inside a block) followed by an OPEN operation
(to remove small isolated noise blobs). The order is significant:
applying OPEN first would erode a block with a glare hole into
disconnected fragments that the subsequent area filter would reject.

\textit{6) Contour extraction and filtering.} External contours
are extracted with \texttt{cv::findContours}. Contours whose
pixel area falls outside configurable bounds (\texttt{min\_block\_area},
\texttt{max\_block\_area}) are discarded.

\textit{7) Centroid computation.} For each surviving contour, the
centroid is computed using image moments: $c_x = M_{10}/M_{00}$,
$c_y = M_{01}/M_{00}$.

\textit{8) Depth lookup.} The depth at the centroid pixel is obtained
by reading a small window of depth pixels around the centroid and
taking the median. The median is robust to noisy zero readings and
single-pixel spike errors that are characteristic of structured-light
and stereo depth sensors.

\textit{9) 3D projection.} The pixel coordinates and depth value are
combined with the camera's intrinsic calibration (obtained via the
\texttt{CameraInfo} topic) to produce a 3D point in the camera frame.
The OpenCV \texttt{PinholeCameraModel::projectPixelTo3dRay} function
produces a unit-length ray from the camera through the pixel, which
is scaled by the measured depth.

\textit{10) TF transformation.} The 3D point in the camera frame is
transformed to \texttt{panda\_link0} via \texttt{tf\_buffer\_.transform}.
A failed transform (typically due to missing or stale TF data) causes
the candidate to be rejected and the next-best candidate to be tried.

\textit{11) Filtering.} The 3D point is rejected if it is non-finite
(NaN or infinity from invalid depth or TF), if it lies inside a
configurable exclusion sphere around the robot base (used to reject
the Panda's status LED), or if it lies inside any of a configurable
list of additional exclusion spheres (used to mask persistent false
positives at known locations).

\textit{12) Temporal smoothing.} The 3D position is fed into an
exponential moving average (EMA) per block. If no previous position
is stored, the current value is accepted as the initial estimate. If
the stored position is older than a timeout, it is dropped and the
current value becomes a fresh estimate. If the jump from the stored
position exceeds a threshold (while the stored position is fresh), the
candidate is rejected as a probable false positive and the caller
tries the next candidate. Otherwise, the new value is blended into the
stored estimate via EMA with a configurable weight $\alpha$.

\subsection{Player Selection}
\label{sec:selection}

The player selection mechanism, implemented in
\texttt{scripts/player\_selection.py}, observes the live block
detections and infers player intent from changes in block presence.

The node operates in two phases. During the \textit{initialization
phase}, the node waits for a minimum number of blocks to be detected
in stable positions for a settling period. Each block's position at
the end of the settling period is recorded as that block's ``slot''.

During the \textit{monitoring phase}, the node subscribes to each new
\texttt{/detected\_blocks} message and checks whether each known
slot's block is still present at (or near) its slot position. If a
block has been continuously absent for a threshold time
(\texttt{missing\_hold\_sec}) and a cooldown has elapsed since the
last selection event, a \texttt{/player\_selection} message is
published with the corresponding block id and color.

If all known blocks disappear simultaneously (typically because the
game ended and the player removed everything), the node resets after
a delay and re-enters the initialization phase. This allows the system
to recover for a new round without restarting the node.

\subsection{Motion Planning}
\label{sec:motion}

Motion is handled by \texttt{newmotion/motion\_moveit\_node.cpp}, which
uses MoveIt for collision-aware joint-space planning on the Panda's
\texttt{arm} planning group.

The node subscribes to \texttt{/target\_sequence} and treats each
incoming message as a complete dispatch: the entire sequence is
executed as one operation. For each block in the sequence, the node
performs a hover-point-hover motion:

\begin{enumerate}
\item Plan and execute a motion to a hover pose above the block at
a configurable height (\texttt{min\_hover\_z}).
\item Plan and execute a downward motion to a lower point pose at
a configurable clearance above the block, indicating the target to
the player.
\item Plan and execute a return motion to the hover pose.
\end{enumerate}

After the last block in the sequence, the arm returns to a home pose
if configured. Throughout the sequence, the node publishes
\texttt{/motion\_status} messages of the form
\texttt{STATE:sequence\_id:block\_id} (e.g.,
\texttt{POINTING:1:2} or \texttt{AT\_TARGET:1:3}) so that the game
node can track progress and distinguish events from the current
sequence from any stragglers from a previous one.

Joint-space planning was chosen over Cartesian path execution after
experimentation: the lab Panda produced smoother and more reliable
motion with planned joint targets than with
\texttt{computeCartesianPath}, which occasionally produced jerky or
infeasible plans for the dip motion in our hover-point-hover pattern.

\subsection{Game State Machine}
\label{sec:game_state}

The \texttt{game\_node} owns the overall game flow as a state machine
(Fig.~\ref{fig:state_machine}). The states are: \texttt{IDLE} (initial),
\texttt{WAITING\_FOR\_BLOCKS} (requires three stable detections before a
round), \texttt{SHOWING\_SEQUENCE} (sequence dispatched to motion),
\texttt{WAITING\_PLAYER} (player reproduces the sequence by removing
blocks), \texttt{CHECKING\_INPUT} (verifying the player's choices),
\texttt{ROUND\_PAUSE} (brief delay before the next round), and the
terminal states \texttt{MOTION\_FAILED} and \texttt{GAME\_OVER}.

Transitions are driven by ROS messages: motion status updates trigger
the move from \texttt{SHOWING\_SEQUENCE} to \texttt{WAITING\_PLAYER},
player selection events trigger the move into \texttt{CHECKING\_INPUT},
and motion failures or sequence-completion-with-error transition into
\texttt{MOTION\_FAILED} or \texttt{GAME\_OVER}. The
\texttt{ROUND\_PAUSE} timer determines how quickly the next round
begins after a successful round.

\begin{figure}[t]
\centering
\fbox{\rule{0pt}{5cm}\rule{0.85\columnwidth}{0pt}}
% \includegraphics[width=\columnwidth]{figures/game_state_machine.png}
\caption{Game node state machine. Transitions are triggered by ROS
messages (motion status, player selection) or internal timeouts
(round pause, player timeout).}
\label{fig:state_machine}
\end{figure}

\section{Design Decisions}
\label{sec:decisions}

Many design choices were made during development. This section
analyzes three decisions that had the largest impact on the system's
behavior and the team's effort budget. Table~\ref{tab:decisions}
summarizes them; the subsections that follow give the full reasoning.
For each, we describe the implemented choice, a credible alternative
path that could have produced a working minimum-viable demonstration,
and the reasons for the chosen path.

\begin{table}[t]
\centering
\caption{Major Design Decisions: Summary}
\label{tab:decisions}
\renewcommand{\arraystretch}{1.2}
\begin{tabular}{p{0.30\columnwidth}p{0.25\columnwidth}p{0.30\columnwidth}}
\toprule
\textbf{Decision} & \textbf{Chose} & \textbf{Alternative} \\
\midrule
Block detection & HSV thresholding & Trained CNN (YOLO) \\
False-positive filtering & 2D image crop & 3D workspace box \\
Motion dispatch & Full sequence per message & One target at a time \\
\bottomrule
\end{tabular}
\end{table}

\subsection{Decision~1: HSV Thresholding vs. Deep Learning Detection}

\textit{Chosen approach.} Block detection uses classical HSV color
thresholding with morphological cleanup, as described in
Section~\ref{sec:vision}. Each color is configured with a manually
tuned HSV range stored in YAML and adjustable without recompilation.

\textit{Alternative considered.} A modern alternative would have been
to train a small convolutional neural network (e.g., YOLOv8-nano or a
purpose-built segmentation model) on a labeled dataset of block images
captured under the lab's actual lighting conditions. The trained
model would produce bounding boxes or instance masks per color, which
could feed into the same downstream depth/projection/smoothing stages.

\textit{Rationale for the chosen path.} HSV was selected for several
reasons. First, the project scope (one semester) did not include
sufficient time to collect a labeled dataset, train, and tune a neural
detector. Second, the lab environment is controlled enough that the
block colors are well-separated in HSV space; manual tuning suffices.
Third, HSV thresholds are inspectable and tunable in YAML without
retraining---when illumination shifts, a team member can adjust four
numbers and observe the result immediately. Fourth, the CPU cost is
negligible, leaving headroom for other processing. The accepted
trade-off is that the system is sensitive to lighting changes and
requires per-deployment tuning, which is acceptable for a fixed-lab
demonstration.

\subsection{Decision~2: Image-Space Crop vs. 3D Workspace Bounding}

\textit{Chosen approach.} False-positive detections from outside the
play area are suppressed by an image-space crop: a configurable
rectangle in pixel coordinates is applied before color detection,
and pixels outside the rectangle are masked to pure black.

\textit{Alternative considered.} A natural alternative is a 3D
workspace bounding box: a region in \texttt{panda\_link0} coordinates
that defines the valid play area, with detections outside the region
filtered out post-projection. This alternative was implemented multiple
times during the project under various names (\texttt{workspace\_*}
parameters, \texttt{target\_workspace\_*} parameters).

\textit{Rationale for the chosen path.} The 3D workspace approach
proved consistently brittle in practice. Its correctness depends on
accurate TF data and stable depth measurements; if either drifted
slightly, the bounds would intermittently reject real blocks, and the
team would respond by widening the bounds until they no longer filtered
anything useful. The cycle of ``add workspace, blocks get rejected,
widen workspace, doesn't filter anymore'' repeated several times in
the commit history.

The image-space crop avoids these failure modes (Fig.~\ref{fig:crop}).
Pixel coordinates do not depend on TF or depth accuracy: the crop is
fixed relative to the camera image, and as long as the camera mount
does not move, the crop remains correct.

\begin{figure}[t]
\centering
\fbox{\rule{0pt}{4cm}\rule{0.85\columnwidth}{0pt}}
% \includegraphics[width=\columnwidth]{figures/crop_demo.png}
\caption{Image-space crop applied as a pre-detection mask. The full
camera frame (left) contains the play area plus background sources
of false positives (operator clothing, robot status LED, ambient
objects). After masking (right), only the play area survives into the
HSV detection stage, eliminating these false-positive sources entirely.}
\label{fig:crop}
\end{figure} Tuning is visual: the team subscribes to a
\texttt{debug\_overlay} topic in \texttt{rqt\_image\_view} that draws
the current crop rectangle as a green outline on the live camera
image, and adjusts the four pixel coordinates until the rectangle
tightly wraps the play area. The crop also acts as a clean
preprocessing boundary: downstream code does not need to handle the
``what if it's outside the workspace'' case.

The accepted trade-off is that the crop must be re-tuned if the
camera mount physically shifts. In practice, the camera in our setup
is fixed, so this is a one-time setup cost.

\subsection{Decision~3: Sequence-Based Motion vs. Single-Target Motion}

\textit{Chosen approach.} The game node dispatches the entire
sequence to the motion node as a single \texttt{BlockSequence} message.
The motion node executes the full sequence atomically and publishes
status updates tagged with the sequence id.

\textit{Alternative considered.} The earlier design (and a common
pattern in similar systems) is to send one target at a time: the
game publishes a single \texttt{Block} target, the motion node
executes a single hover-point-hover, the motion node reports
completion, and the game node sends the next target. The legacy
\texttt{/target\_block} topic in the repository reflects this earlier
design.

\textit{Rationale for the chosen path.} Sequence-based dispatch
simplifies error handling significantly. With single-target dispatch,
the game node had to maintain a per-target state machine: which
target was just sent, are we waiting for completion, did completion
arrive, is the motion status update for the current or a stale target.
With sequence-based dispatch, the motion node owns this state
internally; the game node only needs to know whether the entire
sequence is in progress, succeeded, or failed.

The sequence id embedded in motion status messages also resolves a
common race condition: motion status messages can arrive after a
new sequence has been started, and the game can ignore them by checking
the sequence id. Without an id, the game would need timing-based
heuristics to distinguish current from stale events.

The accepted trade-offs are that the motion node is more complex
(it manages the sequence loop internally), and that any mid-sequence
retry must reuse the original sequence message rather than rebuilding
from current detections; the latter is implemented as a ``frozen
targets'' cache in the game node.

\section{Future Work}
\label{sec:future}

Several aspects of the originally envisioned game were not part of the
demonstrated implementation. Each could be added to the existing
software system with bounded effort.

\textit{Red color reliability.} Reliable red detection in the lab
environment would require either (a) per-session HSV calibration to
adapt to current lighting, (b) depth-aware area filtering that uses
the expected pixel size of a block at the measured depth to reject
larger objects (e.g., a red shirt) even if they pass the HSV threshold,
or (c) replacing color-only detection with fiducial markers
(ArUco/AprilTag) on top of each block, which would also eliminate the
need for per-color tuning altogether at the cost of altering the
physical blocks.

\textit{Difficulty progression.} The game node currently uses a fixed
sequence length (\texttt{base\_length: 3}). The parameter
\texttt{length\_per\_level} already exists and is wired through the
sequence generator; adding a difficulty progression requires
incrementing a level counter after each successful round and using
it as a multiplier. Visual and audio feedback for level transitions
would be added on the UI side.

\textit{Gesture-based player selection.} The current slot-based
selection works reliably but requires the player to physically move
blocks. A gesture-based alternative would let the player point at
the block they wish to select. This could be implemented with
MediaPipe Hands or a similar lightweight pose-estimation model
running as an additional ROS node, publishing to a
\texttt{/player\_selection} topic with the same message format. The
game node would not need any changes.

\textit{Voice and audio feedback.} The browser UI already has
infrastructure for sound playback. Adding text-to-speech announcements
of game state changes, encouragement on correct selections, and
distinct sounds for level transitions would significantly enhance the
HRI feel without changing any of the ROS components.

\textit{Multi-player support.} The current
\texttt{player\_selection.py} tracks a single player. Multi-player
support would involve either (a) splitting the play area into per-player
zones and running one player\_selection instance per zone, or
(b) extending the selection event message with a player id derived from
which slot disappeared. The game node would need a per-player score
table.

\section{Lessons Learned}
\label{sec:lessons}

This section summarizes the major challenges encountered during the
project and reflects on how they would be handled differently in a
future project.

\subsection{Design Challenges}

The largest design lesson was that assumptions about scene cleanliness
made at the planning stage do not survive contact with a real lab
environment. The original mock vision node assumed that ``the camera
sees the blocks and nothing else.'' In practice, the camera sees the
robot's status LED, the operator's clothing, background tables, and
ambient objects whose colors happen to overlap with the configured
block colors.

The system therefore needed multiple defensive layers rather than a
single ``correct'' detection method. The final pipeline applies, in
order: an image crop, HSV thresholding, area filtering, base-exclusion
zones, additional exclusion zones for known false-positive locations,
and temporal smoothing with jump rejection. No single layer is
sufficient, but the layered design tolerates the failure of any one
of them.

In a future project, we would budget explicit time at the start for
identifying ``what does the camera actually see'' in the deployment
environment, rather than assuming the scene matches the design's
ideal.

\subsection{Implementation Challenges}

The most consistently difficult implementation problem was
synchronization between the depth and color images from the RealSense
camera. The two image streams arrive with separate timestamps that do
not always align perfectly; without explicit handling, the depth
lookup at a centroid pixel could read from a frame whose content no
longer matches the color frame's content.

The solution was a buffered closest-match strategy: depth frames are
held in a small ring buffer, and for each color frame the buffer is
searched for the timestamp-nearest depth frame within a configurable
tolerance. If no frame is within tolerance, the most recent depth is
used as a fallback with a warning.

Beyond synchronization, the morphological cleanup order
(close-then-open) and the HSV range tuning each required several
iterations against real lab data, not synthetic test images. The order
of morphology operations in particular was a small but high-impact
choice: applying open-then-close instead of close-then-open caused
blocks with internal glare holes to be fragmented and subsequently
rejected by the area filter.

\subsection{Organizational Challenges}

The most painful organizational lesson was the cost of silent reverts
and uncoordinated changes. The project history includes commits where
substantial recent work was removed (often as a side effect of a
broader rollback intended to restore an earlier runtime behavior). In
each case, the change was technically reasonable in isolation, but the
team member making it did not know about the newer work that was
caught in the rollback.

A specific example: a commit titled ``Restore [earlier] runtime
behavior'' rolled back several files to a previous state. Among the
changes was the removal of a recently added image-cropping feature
that was the next day's planned demonstration improvement. The
rollback was performed because of a separate motion issue; the
person performing it did not see the vision changes as relevant. The
fix required a careful manual merge: keep the rollback's good
defaults, restore the vision feature on top, preserve the
documentation commit added in the interim.

In future projects, we would adopt a stricter coordination protocol:
larger reverts get a heads-up message in the team chat naming the
files they will touch, smaller atomic commits per logical change, and
clear ownership of each major subsystem so the team member responsible
is consulted before changes to ``their'' code.

\subsection{Tooling and Workflow}

A general lesson is that runtime configuration (parameters in
\texttt{game\_params.yaml}) ended up being touched far more often
than any code file. Many bugs and behavior issues were resolved by
parameter changes rather than code changes. In retrospect, the team
should have invested earlier in a configuration management strategy:
clear documentation of what each parameter does, version control for
parameter changes with the same rigor as code, and ideally a runtime
parameter inspection tool that surfaced the active values at start-up.

\section{Conclusion}
\label{sec:conclusion}

We have presented the design and implementation of a memory game
system using a Franka Emika Panda arm and a RealSense RGB-D camera.
The system is composed of independent ROS nodes for vision, player
selection, game logic, and motion, communicating through standard ROS
topics with a small set of custom message types. The vision pipeline
uses classical HSV thresholding with multiple defensive filtering
layers, the player selection mechanism detects block disappearance
from initialization-time slots, and the motion controller uses MoveIt
to execute sequence-based hover-point-hover trajectories.

The project demonstrates that classical computer vision and standard
motion planning---when carefully integrated and tuned against
real-world conditions---can produce a reliable interactive robot
system within a semester-length project. The major design decisions
favored simplicity, tunability, and inspectability over modern but
heavier alternatives; the trade-offs accepted are realistic for a
fixed-deployment educational demonstrator, and the failure modes of
each design choice are well understood.

\begin{thebibliography}{00}
\bibitem{ros} M. Quigley et al., ``ROS: an open-source Robot Operating
System,'' in \textit{ICRA Workshop on Open Source Software}, 2009.
\bibitem{opencv} G. Bradski, ``The OpenCV Library,'' \textit{Dr. Dobb's
Journal of Software Tools}, 2000.
\bibitem{moveit} D. Coleman, I. Sucan, S. Chitta, and N. Correll,
``Reducing the barrier to entry of complex robotic software: a
MoveIt! case study,'' \textit{Journal of Software Engineering for
Robotics}, vol.~5, no.~1, pp.~3--16, 2014.
\bibitem{realsense} L. Keselman, J. Iselin Woodfill, A. Grunnet-Jepsen,
and A. Bhowmik, ``Intel RealSense Stereoscopic Depth Cameras,'' in
\textit{IEEE Conference on Computer Vision and Pattern Recognition
Workshops}, 2017.
\bibitem{franka} T. Haddadin et al., ``The Franka Emika Robot: A
Reference Platform for Robotics Research and Education,''
\textit{IEEE Robotics \& Automation Magazine}, vol.~29, no.~2,
pp.~46--64, 2022.
\end{thebibliography}

\end{document}
